\section{Описание}

Требуется написать реализацию алгоритма поразрядной сортировки.

Основная идея поразрядной сортировки заключается в том, чтобы каждый входной элемент $x$ разбить на разряды и сортировать стабильной сортировкой по каждому разряду.

\pagebreak

\section{Исходный код}

На каждой непустой строке входного файла располагается пара \enquote{ключ-значение}, поэтому создадим новую структуру $Pair$, в которой будем хранить исходную строку и преобразованный в \texttt{uint64\_t} ключ.

Такое решение было принято из-за нехватки памяти, так как существует три варианта хранения номера телефона. Можно хранить ключ как строку и сортировать по строке, но это занимает достаточно много памяти. Можно распарсить строку на три целочисленные переменные, однако в этом случае всё равно потребуется объединять их в одно число для более удобной реализации сортировки.

Поэтому был выбран третий вариант: хранить исходную строку, из которой можно извлечь ключ и значение, а также хранить преобразованный в \texttt{uint64\_t} ключ.

Так как заранее неизвестно количество входных строк, необходимо реализовать динамический массив $Vector$. Для него реализуется правило пяти (деструктор, конструктор копирования, оператор копирования, конструктор перемещения, оператор перемещения), а также дополнительные функции: PushBack, Size и вывод в поток.

Далее необходимо реализовать саму поразрядную сортировку \texttt{RadixSort}. Для неё требуется быстрая стабильная сортировка — в данной работе используется сортировка подсчётом для разряда числа \texttt{CountingSortStep}.

Так как ключ хранится в типе данных \texttt{uint64\_t}, он занимает 8 байт, что соответствует 8 разрядам. Для каждого разряда выделяется дополнительный массив размером 256, так как в одном байте содержится 8 бит, а $2^8 = 256$.

Индекс массива, в который необходимо увеличить значение счётчика, получается побитовым сдвигом ключа на нужное количество бит и побитовой операцией И с маской \texttt{0xFF}, чтобы получить последние 8 бит.

Остальная часть алгоритма представляет собой стандартную реализацию сортировки подсчётом.

Интересным моментом является использование статического массива для хранения дополнительного массива, чтобы не выделять память на каждом шаге алгоритма. Однако после каждой итерации этот массив необходимо очищать, чтобы данные из разных итераций не пересекались.

Также используется двойной буфер для сортировки: вместо выделения нового массива на каждом шаге просто меняются указатели на массивы.

\pagebreak

\begin{table}[!ht]
\centering
\begin{tabular}{|p{7.5cm}|p{7.5cm}|}
\hline
\rowcolor{lightgray}
\multicolumn{2}{|c|}{main.cpp} \\
\hline
\texttt{void CountingSortStep(Vector<Pair>\& input, Vector<Pair>\& output, int shift)} & Функция сортировки подсчётом \\
\hline
\rowcolor{lightgray}
\multicolumn{2}{|c|}{main.cpp} \\
\hline
\texttt{void RadixSort(Vector<Pair>\& pairs)} & Функция поразрядной сортировки \\
\hline
\end{tabular}
\end{table}

\begin{lstlisting}[language=C++]
class Vector {
private:
T* data;
size_t size;
size_t capacity;

void reallocate(size_t new_capacity);

public:
Vector();

Vector(size_t size, T value = T{});

Vector(const Vector& other);

Vector& operator=(const Vector& other);

Vector(Vector&& other) noexcept;

Vector& operator=(Vector&& other) noexcept;

~Vector() noexcept;

size_t Size() const;

void PushBack(T&& value);

friend std::ostream& operator<<(std::ostream& os, const Vector<T>& vector);

T& operator[](size_t idx);

const T& operator[](size_t idx) const;

};

struct Pair {
std::string raw_str;
uint64_t numeric_key;

Pair();

explicit Pair(const std::string& str);

Pair(const Pair& other);

Pair& operator=(const Pair& other);

Pair(Pair&& other) noexcept;

Pair& operator=(Pair&& other) noexcept;

~Pair();

friend std::ostream& operator<<(std::ostream& os, const Pair& pair);
};
\end{lstlisting}

\pagebreak

\section{Консоль}

\begin{alltt}
root@74b1f15b9ece:/workspaces/Discrete-Analysis-MAI# cmake -B build
root@74b1f15b9ece:/workspaces/Discrete-Analysis-MAI# cmake --build build
root@74b1f15b9ece:/workspaces/Discrete-Analysis-MAI# ./build/lab1/lab1
+7-495-1123212    xGfxrxGGxrxMMMMfrrrG
+375-123-1234567  xGfxrxGGxrxMMMMfrrr
+7-495-1123212    xGfxrxGGxrxMMMMfrr
+375-123-1234567  xGfxrxGGxrxMMMMfr
+7-495-1123212    xGfxrxGGxrxMMMMfrrrG
+7-495-1123212    xGfxrxGGxrxMMMMfrr
+375-123-1234567  xGfxrxGGxrxMMMMfrrr
+375-123-1234567  xGfxrxGGxrxMMMMfr
\end{alltt}

\pagebreak
